\documentclass[10pt,a4paper]{amsart}
\usepackage[margin=19mm]{geometry}
\usepackage{amsmath,amssymb,mathtools}
\usepackage[scr=rsfs,cal=euler]{mathalfa}
\usepackage{xfrac}
\usepackage[unicode,hidelinks,bookmarks=false]{hyperref}
\newcommand{\ie}{i.e.\ }
\newcommand{\cf}{cf.\ }
\newcommand{\resp}{respectively\ }
\newcommand{\Subsection}{\subsection}
\providecommand{\nobreakdash}{\nobreak\hbox{-}}
\setlength{\emergencystretch}{2em}
\multlinegap0pt
\usepackage{mathtools,tikz-cd}
\usepackage{algorithm}
\usepackage{algpseudocode}
\usepackage{color}
\def\red#1{\textcolor{red}{#1}} 
\def\blue#1{\textcolor{blue}{#1}} 
\usepackage{ulem}
\newcommand{\Z}{\mathbb{Z}}
\newcommand{\R}{\mathbb{R}}
\newcommand{\cF}{\mathcal{F}}
\newcommand{\cV}{\mathcal{V}}
\newcommand{\cB}{\mathcal{B}}
\newcommand{\cI}{\mathcal{I}}
\newcommand{\cU}{\mathcal{U}}
\newcommand{\cP}{\mathcal{P}}
\DeclareMathOperator{\GL}{GL}
\DeclareMathOperator{\im}{im}
\DeclareMathOperator{\id}{Id}
\DeclareMathOperator{\Span}{span}
\DeclareMathOperator{\SAT}{SAT}
\DeclareMathOperator{\UNSAT}{UNSAT}
\DeclareMathOperator{\supp}{supp}
\newtheorem{theorem}{Theorem}[section]
\newtheorem{proposition}[theorem]{Proposition}
\newtheorem{lemma}[theorem]{Lemma}
\newtheorem{corollary}[theorem]{Corollary}
\newtheorem{definition}[theorem]{Definition}
\newtheorem{example}[theorem]{Example}
\newtheorem{remark}[theorem]{Remark}
\begin{document}
\title[On the toric lifting properties for simplicial 3-spheres]{On the toric lifting properties for simplicial $3$-spheres}
\author{Suyoung Choi, Yuanxin Guan, Hyeontae Jang, Zhi Lü}
\date{}
\maketitle
\noindent\emph{Source and licence.} On the toric lifting properties for simplicial $3$-spheres. Version arXiv:2607.18706v1 (CC BY 4.0). This material is distributed under the Creative Commons Attribution 4.0 International licence, http://creativecommons.org/licenses/by/4.0/, as recorded in the arXiv OAI licence metadata for this exact version and shown on the abstract page https://arxiv.org/abs/2607.18706v1. Full rights notice: licenses/arxiv-2607.18706-CC-BY-4.0.txt.\par\smallskip
\noindent\textit{Editorial scope.} Counted unit: the original \texttt{theorem} environment \texttt{thm:main} at source lines 273--276 together with its complete proof at lines 277--304; the delivered document keeps source lines 145--305 so that every referenced label and every citation resolves inside it. Native editable source is the paper's own concatenated TeX; the AMS wrapper is editorial formatting only. No publisher class, upstream script or unrelated artwork is redistributed.
\section{Universal Complexes} \label{universal cpx}
Let $A=(a_1,\dots,a_{15})$ be the $4\times 15$ $\Z_2$-matrix whose columns are the nonzero vectors of $\Z_2^4$ listed in lexicographic order
$$
    A= \left(
    \begin{array}{ccccccccccccccc}
        0&0&0&0&0&0&0&1&1&1&1&1&1&1&1\\
        0&0&0&1&1&1&1&0&0&0&0&1&1&1&1\\
        0&1&1&0&0&1&1&0&0&1&1&0&0&1&1\\
        1&0&1&0&1&0&1&0&1&0&1&0&1&0&1
    \end{array} \right),
$$
and let $X\bigl(\Z_2^4\bigr)$ be the corresponding universal complex formed by all linear independent subsets of $\Z_2^4$.
Thus the vertex set of~$X(\Z_2^4)$ is $\cV\coloneq\cV(X(\Z_2^4)) = \{a_1,\dots,a_{15}\}$, and its facet set is
$$
    \cF\coloneq\cF(X(\Z_2^4)) =\Bigl\{ \{a_{i_1},a_{i_2},a_{i_3},a_{i_4}\}\subset \cV \Bigm| \det(a_{i_1},a_{i_2},a_{i_3},a_{i_4})\equiv 1 (\bmod 2) \Bigr\}.
$$
Then $|\cF|=\frac{|\GL(4,\Z_2)|}{4!}=\frac{(2^4-1)(2^4-2)(2^4-2^2)(2^4-2^3)}{4!}=840$.

With our choice of ordering, and viewing each $a_i$ as an integer vector in $\{0,1\}^4\subset \Z^4$, the set of facets of $X(\Z_2^4)$ whose determinants have absolute value $3$ is
$$
    B= \left\{
        \begin{aligned}
            &\{a_3,a_5,a_9,a_{14}\},\ \{a_3,a_6,a_{10},a_{13}\},\ \{a_5,a_6,a_{11},a_{12}\},\\
            &\{a_7,a_9,a_{10},a_{12}\},\ \{a_7,a_{11},a_{13},a_{14}\}
        \end{aligned}
    \right\}.
$$
We call $B$ the \emph{standard bad block}.

Note that the action of $\GL(4,\Z_2)$ on $X(\Z_2^4)$ permutes the facets of $X(\Z_2^4)$.
One interesting fact, originally established in the proof of \cite[Theorem~6.1]{Choi-Jang-Vallee2025JLMS}, is that $\cF$ admits a partition
$$
    \cF=\bigsqcup_{\alpha=1}^{168} B_\alpha
$$
into $168$ subsets of size $5$, and for each $\alpha$ there exists $g_\alpha\in \GL(4,\Z_2)$ such that
$$
    g_\alpha(B_\alpha)=B.
$$
We may assume that $B_1=B$.

For a facet $\tau=\{u_1,u_2,u_3,u_4\}\in \cF$, set
$$
    P(\tau)\coloneq \{\tau\}\cup \bigl\{ \{u_i\}\cup \{u_j+u_i\mid j\neq i\} \bigm| i=1,2,3,4 \bigr\}.
$$
With this notation, the standard bad block is
$$
    B=P(\{a_3,a_5,a_9,a_{14}\}).
$$
The distinct sets $P(\tau)$, as $\tau$ runs over the facets of $X(\Z_2^4)$, form the bad-block partition.
In particular, for every basis $\{u,p,q,r\}$ of $\Z_2^4$, the five facets
\begin{align*}
    &\{u,p,q,r\},\quad \{u,p+u,q+u,r+u\},\quad  \{u+p,p,q+p,r+p\}, \\
    &\{u+q,p+q,q,r+q\},\quad \{u+r,p+r,q+r,r\}  
\end{align*}
belong to the same block of the bad-block partition.

The following simple observation explains why bad blocks control the validity of integral lifts.

\begin{lemma}\label{lem:avoid-standard-bad-block}
    Let $Y$ be a pure $3$-dimensional simplicial complex, and let $\eta^\R$ be a mod~$2$ characteristic map over~$Y$.
    Suppose that no facet image of $\eta^\R$ lies in the standard bad block~$B$.
    Then the map obtained by taking the $\{0,1\}$-representative of each vector $\eta^\R(v)\in \Z_2^4$ is a $\{0,1\}$-lift of $\eta^\R$.
\end{lemma}
\begin{proof}
    Let $\eta$ be the $\{0,1\}$-representative lift of $\eta^\R$.
    For each facet $\sigma$ of $Y$, the image~$\widehat{\eta^\R}(\sigma)$ is a facet of $X(\Z_2^4)$.
    The determinant of the corresponding $0$-$1$ matrix is odd.
    Since every $4\times4$ 0-1 matrix has determinant of absolute value at most~$3$, an odd determinant is equal to~$\pm 1$ or~$\pm 3$.
    By the definition of $B$, the only facets of $X(\Z_2^4)$ whose $0$-$1$ determinant has absolute value $3$ are precisely the facets in~$B$.
    Since $\widehat{\eta^\R}(\sigma)\notin B$, this determinant has absolute value $1$.
    Thus the images of the vertices of every facet of $Y$ form a unimodular basis of $\Z^4$, and so $\eta$ is a $\{0,1\}$-lift.
\end{proof}

We note that a pure $3$-dimensional simplicial complex $Y$ supports a characteristic map $\lambda^\R$ if and only if there is a nondegenerate simplicial map $\widehat{\lambda^\R} \colon Y \to X(\Z_2^4)$.
For a facet $\sigma = \{i_1, i_2, i_3, i_4\}$ of $Y$, we write
$$
    \widehat{\lambda^\R}(\sigma) \coloneq\{\widehat{\lambda^\R}(i_1),\widehat{\lambda^\R}(i_2),\widehat{\lambda^\R}(i_3),\widehat{\lambda^\R}(i_4)\}\in\cF,
$$
and define
$$
    \Sigma(\lambda^\R)\coloneq \Bigl\{\widehat{\lambda^\R}(\sigma)\ \Bigm|\ \sigma\in\cF(Y)\Bigr\}.
$$

\begin{lemma}[Bad-block avoidance criterion]\label{lem:block-avoidance}
    Let $Y$ be a pure $3$-dimensional simplicial complex, and let $\lambda^\R$ be a mod~$2$ characteristic map over $Y$.
    If $\Sigma(\lambda^\R)$ is disjoint from some block~$B_\alpha$ of the bad-block partition
    $$
        \cF(X(\Z_2^4))=\bigsqcup_{\alpha=1}^{168}B_\alpha,
    $$
    then $\lambda^\R$ admits a lift.
    Or equivalently, if $\Sigma(\lambda^\R)$ meets fewer than $168$ blocks of this partition, then $\lambda^\R$ admits a lift.
\end{lemma}
\begin{proof}
    Choose $\alpha$ such that $\Sigma(\lambda^\R)\cap B_\alpha=\varnothing$.
    Since $g_\alpha(B_\alpha)=B$, we have
    $$
        \Sigma(g_\alpha\lambda^\R)\cap B=\varnothing.
    $$
    Let $\widetilde\lambda$ be the map obtained by taking the $\{0,1\}$-representative of each vector~$g_\alpha\lambda^\R(v)$.
    By Lemma~\ref{lem:avoid-standard-bad-block}, $\widetilde\lambda$ is a $\{0,1\}$-lift of $g_\alpha\lambda^\R$.
    
    The reduction map $\GL(4,\Z) \to \GL(4,\Z_2)$ is surjective, since $\GL(4,\Z_2)$ is generated by elementary matrices, each of which has an integral unimodular lift.
    Choose an integral lift $G_\alpha\in \GL(4,\Z)$ of $g_\alpha$, and define
    $$
        \lambda\coloneq G_\alpha^{-1}\widetilde\lambda.
    $$
    Then $\lambda$ is an integral characteristic map and
    $$
        \lambda=G_\alpha^{-1}\widetilde\lambda \equiv g_\alpha^{-1}(g_\alpha\lambda^\R)=\lambda^\R \pmod 2.
    $$
    Hence $\lambda$ is a lift of $\lambda^\R$.
\end{proof}

\begin{proposition}\label{prop:168}
    Let $Y$ be a pure $3$-dimensional simplicial complex, and let $\lambda^\R$ be a mod~$2$ characteristic map over $Y$.
    If $|\Sigma(\lambda^\R)|<168$, then $\lambda^\R$ admits a lift.
    In particular, if $K$ is a simplicial $3$-sphere with $f_3(K)<168$, then any mod~$2$ characteristic map over $K$ admits a lift.
\end{proposition}
\begin{proof}
    Since the bad-block partition has $168$ blocks and $|\Sigma(\lambda^\R)|<168$, the set~$\Sigma(\lambda^\R)$ is disjoint from at least one block.
    The result follows from Lemma~\ref{lem:block-avoidance}.
    The final assertion follows because $|\Sigma(\lambda^\R)|\le f_3(K)$.
\end{proof}


\section{Toric lifting properties for 3-spheres} \label{toric lifting property}

Based upon the bad-block avoidance criterion established in Section~\ref{universal cpx}, we first prove our  vertex bound theorem for general simplicial 3-spheres and derive a corollary for neighborly  simplicial 3-spheres.

\begin{theorem}\label{thm:main}
    Let $K$ be a simplicial $3$-sphere with at most $20$ vertices.
    Then $K$ has the toric lifting property.
\end{theorem}

\begin{proof}
    Let $K$ be a simplicial $3$-sphere on $[m]$ with $m \leq 20$.
    Then, by the \emph{upper bound theorem}~\cite{Stanley1975}, the number $f_3(K)$ of facets is at most $\frac{m(m-3)}{2}$.
    In particular, if $m \leq 19$, then $f_3(K) \leq \frac{19 \cdot 16}{2} = 152$.
    Now suppose that $m=20$ and $K$ supports a mod~$2$ characteristic map, and let $\widehat{\lambda^\R} \colon K\to X\bigl(\Z_2^4\bigr)$ be the corresponding map.
    Since $\widehat{\lambda^\R}$ is nondegenerate, the endpoints of every edge of $K$ must be sent to distinct vertices of $X(\Z_2^4)$, and, hence, the $1$-skeleton of $K$ is $15$-colorable.
    
    Let $n_1,\dots,n_{15}$ be the color class sizes. 
    Then $n_1+\cdots+n_{15}=20$, and the number of edges satisfies
    $$
        f_1(K) \leq \binom{20}{2}-\sum_{i=1}^{15}\binom{n_i}{2}.
    $$
    Since $20$ vertices are distributed among $15$ classes, at least five pairs of vertices must lie in the same color class, and therefore
    $$
        \sum_{i=1}^{15}\binom{n_i}{2}\ge 5.
    $$
    Thus $f_1(K)\le \binom{20}{2}-5=185$.
    
    Now $K$ is a simplicial $3$-sphere, so every triangle lies in exactly two tetrahedra. 
    Hence $2f_2(K)=4f_3(K)$, that is, $f_2(K)=2f_3(K)$.
    Applying Euler's relation, we obtain
    $$
        20-f_1(K)+2f_3(K)-f_3(K)=0,
    $$
    and, hence, $f_3(K)=f_1(K)-20\le 185-20=165$.
    
    Hence, by Proposition~\ref{prop:168}, $K$ has the toric lifting property.
\end{proof}


\begin{thebibliography}{99}
\bibitem[Choi-Jang-Vallee2025JLMS]{Choi-Jang-Vallee2025JLMS} Suyoung Choi, Hyeontae Jang, and Mathieu Vall\'ee, \emph{Toric wedge induction and toric lifting property for piecewise linear spheres with few vertices}, J. Lond. Math. Soc. (2) \textbf{112} (2025), no.~1, Paper No. e70231, 24. \MR{4931929}
\bibitem[Stanley1975]{Stanley1975} Richard~P. Stanley, \emph{The upper bound conjecture and {C}ohen-{M}acaulay rings}, Studies in Appl. Math. \textbf{54} (1975), no.~2, 135--142. \MR{458437}
\end{thebibliography}
\end{document}
